\customSection[sec:design]{Обґрунтування та проєктування вдосконаленої системи}{0em}

Сформульовані вимоги визначають напрями розвитку системи Dailu.
У цьому розділі коротко описано наявну архітектурну основу, на яку спираються нововведення,
після чого обґрунтовано й спроєктовано дві нові складові — AI-підсистему та IoT-канал збору
активності, а також уточнено технологічний стек.

\customSubSection[sec:design_base]{Наявна архітектурна основа}{0em}

Серверну частину Dailu побудовано як модульний моноліт за принципами
предметно-орієнтованого проєктування та чистої архітектури.
Систему поділено на обмежені контексти (\textit{Habit}, \textit{HabitEntry},
\textit{HabitUser}, \textit{Identity}, \textit{Tag}), кожен з яких відповідає за самодостатню
частину предметної області та взаємодіє з іншими лише через визначені інтерфейси.
Сценарії застосування реалізовано за принципом розділення команд і запитів (CQRS), а взаємодія
між модулями відбувається через інтеграційні події, а не прямі виклики.

Ця подієва основа має принципове значення для дальших нововведень.
Автоматичний збір активності вже побудовано на публікації події
\textit{IntegrationActivitiesDetected}, на яку підписаний модуль записів.
Завдяки цьому додавання нового джерела активності (у нашому випадку — IoT-пристрою) не потребує
зміни логіки формування записів: достатньо, щоб нове джерело публікувало ту саму подію.
Так само слабке зв'язування модулів дає змогу вбудувати AI-підсистему як окремий модуль, не
порушуючи наявних меж відповідальності.

\customSubSection[sec:design_ai]{Обґрунтування підходу до AI-підсистеми}{1em}

Інтелектуальні функції системи можна реалізувати двома шляхами: навчанням і розгортанням
власних моделей машинного навчання або використанням готових великих мовних моделей через API
зовнішнього провайдера.
Для навчального проєкту з однією командою розробників і обмеженим обсягом розмічених даних
власне навчання моделей є невиправдано складним і не гарантує якості.
Тому за основу обрано підхід із використанням готової великої мовної моделі, доповнений подачею
до неї контексту даних користувача за принципом генерації, доповненої вибіркою
(retrieval-augmented generation, RAG), — це забезпечує прийнятну якість висновків без потреби у
власному навчальному конвеєрі та великому наборі розмічених даних.

AI-підсистему спроєктовано як окремий обмежений контекст, що взаємодіє з рештою системи через
визначені інтерфейси.
Вона охоплює три функції відповідно до вимог:

\begin{itemize}[noitemsep, topsep=2pt, leftmargin=*, label=--]
    \item \textbf{рекомендації та інсайти} — підсистема отримує агреговану історію виконання звичок користувача (серії, частоту, пропуски), формує на її основі структурований запит до мовної моделі й повертає стислі персоналізовані поради;
    \item \textbf{автоматична категоризація} — під час створення звички її назва й опис передаються моделі, яка пропонує доречні теги; користувач може прийняти або скоригувати пропозицію, зберігаючи контроль над даними;
    \item \textbf{діалоговий асистент} — користувач спілкується з мотиваційним асистентом природною мовою; до запиту додається контекст його звичок, що робить відповіді релевантними.
\end{itemize}

Ключові архітектурні рішення тут — ізоляція звернень до зовнішнього AI-провайдера за окремим
інтерфейсом (що дає змогу змінювати провайдера без впливу на решту коду), стійкість до
тимчасових збоїв зовнішнього сервісу та контроль користувача над обсягом даних, які передаються
назовні, відповідно до вимоги приватності.

\customSubSection[sec:design_iot]{Обґрунтування та проєктування IoT-складової}{1em}

Мета IoT-складової — охопити автоматизованим збором ті звички, що не залишають цифрового сліду
в програмних сервісах.
Як показано в огляді, для таких випадків ефективним є активний збір даних простим фізичним
пристроєм \citecustom{7}.
Тому як демонстраційний пристрій обрано мінімалістичний варіант — фізичну кнопку-трекер на базі
мікроконтролера ESP32 з вбудованим модулем Wi-Fi (умовна назва «Dailu Button»).

Принцип роботи полягає в такому: натискання кнопки відповідає одному факту виконання прив'язаної
до пристрою бінарної звички.
Пристрій установлює захищене з'єднання із серверною частиною та надсилає подію про натискання;
сервер перетворює її на подію активності й далі опрацьовує тим самим механізмом, що й активність
GitHub чи Strava.
Такий вибір обґрунтовано трьома міркуваннями:

\begin{enumerate}[noitemsep, topsep=2pt, leftmargin=*, label=\arabic*)]
    \item простота реалізації — ESP32 є доступною й добре задокументованою платформою, а сценарій з однією кнопкою мінімізує апаратну складність і дає змогу зосередитися на програмній інтеграції;
    \item архітектурна узгодженість — пристрій виступає ще одним джерелом активності й вписується в наявну подієву модель без зміни логіки формування записів;
    \item демонстративність — фізична дія наочно ілюструє ідею переходу від ручного обліку до автоматизованого збору з реального світу.
\end{enumerate}

Для безпечної взаємодії кожен пристрій отримує унікальний ключ, що прив'язується до облікового
запису користувача та звички.
Обмін даними відбувається за захищеним протоколом; на боці сервера приймання подій пристрою
реалізовано окремим ендпоінтом, який після перевірки ключа публікує вже знайому системі подію
\textit{IntegrationActivitiesDetected} із відповідним джерелом.
Захист від дублювання записів забезпечується наявним механізмом перевірки за зовнішнім
ідентифікатором події.
Спрощений скетч прошивки, що надсилає подію натискання на сервер, наведено в лістингу~\ref{lst:iot}.

\begin{lstlisting}[style=nocodeframe, language=ArduinoCpp, caption={Надсилання події натискання кнопки з пристрою ESP32}, label={lst:iot}]
void loop() {
    if (digitalRead(BUTTON_PIN) == LOW) {
        HTTPClient http;
        http.begin(client, SERVER_URL);
        http.addHeader("Content-Type", "application/json");
        http.addHeader("X-Device-Key", DEVICE_KEY);

        String body = "{\"pressedAt\":\"" + isoTimestamp() + "\"}";
        int code = http.POST(body);

        if (code == 200) blink(LED_PIN);
        http.end();
        delay(DEBOUNCE_MS);
    }
}
\end{lstlisting}

\customSubSection[sec:design_stack]{Уточнений технологічний стек}{1em}

Технологічну основу успадковано від наявної системи й доповнено засобами для нових складових.
Серверну частину реалізовано на платформі .NET мовою C\#, доступ до даних — через Entity
Framework Core з провайдером Npgsql, сховище — PostgreSQL.
Клієнтську частину побудовано як односторінковий застосунок на Angular.
Нові складові додають до стеку інтеграцію із зовнішнім AI-провайдером через його API на боці
сервера та прошивку IoT-пристрою мовою C++ у середовищі для ESP32.
Стислий перелік технологій наведено в \tabref{tab:stack}.

\begin{center}
    \begin{xltabular}{\textwidth}{| >{\raggedright\arraybackslash}X | >{\raggedright\arraybackslash}X |}
        \caption{Технологічний стек удосконаленої системи} \label{tab:stack} \\
        \hline
        \textbf{Складова} & \textbf{Технологія} \\ \hline
        \endfirsthead

        \hline
        \textbf{Складова} & \textbf{Технологія} \\ \hline
        \endhead

        Платформа сервера & .NET, C\# \\ \hline
        Вебфреймворк & ASP.NET Core \\ \hline
        Доступ до даних & Entity Framework Core, Npgsql \\ \hline
        База даних & PostgreSQL \\ \hline
        Клієнтський застосунок & Angular \\ \hline
        AI-підсистема & Інтеграція з LLM-провайдером через API \\ \hline
        IoT-пристрій & ESP32 (C++), Wi-Fi, HTTPS \\ \hline
    \end{xltabular}
\end{center}

Узагальнену архітектуру вдосконаленої системи — з успадкованими обмеженими контекстами, новим
модулем інтелектуального аналізу \textit{AI} та каналом збору активності від IoT-пристрою —
показано на \textbf{рис.~\ref{fig:modules}}.

\setfigure[fig:modules]{0.7\textwidth}{assets/diagram-modules}{Архітектура вдосконаленої системи: модулі серверної частини та зовнішні залежності}

Описані архітектурні рішення повністю покривають сформульовані у першому розділі вимоги.
Вони спираються на наявну модульну основу й розширюють її двома новими каналами — інтелектуальним
аналізом даних та збором фізичної активності, — не порушуючи меж відповідальності модулів.
Детальний розгляд реалізації цих складових виноситься в наступний розділ роботи.
